\documentclass[../DoAn.tex]{subfiles}
\begin{document}

\begin{center}
    \Large{\textbf{TÓM TẮT NỘI DUNG ĐỒ ÁN}}\\
\end{center}
\vspace{1cm}
Người dân và doanh nghiệp Việt Nam thường xuyên cần tra cứu pháp luật, nhưng hệ thống văn bản quy phạm pháp luật có khối lượng lớn, thay đổi liên tục và được viết bằng ngôn ngữ chuyên môn. Các công cụ tra cứu hiện có chủ yếu tìm kiếm theo từ khoá nên đòi hỏi người dùng phải biết trước thuật ngữ pháp lý, trong khi các chatbot dựa trên mô hình ngôn ngữ lớn có thể trả lời trôi chảy nhưng dễ bịa đặt điều luật hoặc dẫn văn bản đã hết hiệu lực. Hướng sinh văn bản tăng cường truy xuất (RAG) khắc phục một phần hạn chế này bằng cách cho mô hình đọc tài liệu trước khi trả lời, song một bước truy xuất cố định thường bỏ sót các điều luật liên quan hoặc các quy định mới.

Đồ án lựa chọn hướng kết hợp RAG, đồ thị tri thức và tác tử AI trên Bộ Pháp điển do Bộ Tư pháp biên soạn, vì đây là nguồn dữ liệu chính thống đã được hệ thống hoá theo chủ đề và có sẵn quan hệ giữa các điều luật. Hệ thống được xây dựng gồm một tầng tiền xử lý phân loại yêu cầu, một tác tử gọi công cụ có thể tự quyết định tìm kiếm lai (vector kết hợp BM25, tự kèm điều được dẫn chiếu và khoản ngoại lệ từ đồ thị điều luật), đọc toàn văn điều luật và kiểm tra quy định hiện hành trên web, cùng bộ nhớ dài hạn có kiểm soát quyền riêng tư. Ngoài hỏi đáp, hệ thống cho phép người dùng tải lên hợp đồng để được tóm tắt và cảnh báo điều khoản có nguy cơ trái luật; mỗi cảnh báo chỉ được đưa ra khi có căn cứ là điều luật thực sự tra được, kèm đoạn trích nguyên văn đã được kiểm chứng.

Đóng góp chính của đồ án gồm phương pháp truy xuất lai cho tiếng Việt pháp lý, đồ thị dẫn chiếu với kiểm tra ngoại lệ, cơ chế kiểm chứng căn cứ cho tính năng rà soát văn bản, bộ nhớ dài hạn bảo vệ dữ liệu nhạy cảm và công cụ quản trị dữ liệu văn bản có theo dõi hiệu lực. Trên 50 câu hỏi chấm mù bằng mô hình ngôn ngữ, truy xuất lai nâng P@10 từ 0,368 lên 0,544 và MRR từ 0,280 lên 0,471 so với cách ban đầu, dù chưa tốt hơn BM25 đơn lẻ một cách có ý nghĩa thống kê; tác tử GPT-4.1 gọi đủ công cụ bắt buộc ở 97,6\% trường hợp và bộ nhớ dài hạn đạt cả 14 kịch bản mà không lộ dữ liệu nhạy cảm.

\begin{flushright}
Sinh viên thực hiện\\
\begin{tabular}{@{}c@{}}
\textit{(Ký và ghi rõ họ tên)}\\[2cm]
Nguyễn Đức Hải
\end{tabular}
\end{flushright}

\end{document}
